\documentclass[11pt]{article}
\usepackage{amsmath}
\usepackage{amssymb}
\usepackage{amsthm}
\usepackage{enumitem}
\usepackage[margin=1in]{geometry}
\usepackage{multicol}
\usepackage{tikz}

% Load hyperref to automatically generate the PDF outline from structural tags
\usepackage[pdftex, bookmarks=true, bookmarksnumbered=true]{hyperref}

% Optional: Use the bookmark package for faster rendering and advanced control
\usepackage{bookmark} 

\newenvironment{case}[2]{\par\textbf{Case #1: #2.}}{}
\title{Homework \#2\\{\large Math 104: Intro to Analysis \\ Summer 2026} \\ {\normalsize UC Berkeley}}
\date{July 10, 2026}
\author{Donald Aingworth IV}

\begin{document}
\renewcommand\qedsymbol{OE$\Delta$}
\newcommand{\abs}[1]{\left| #1 \right|}
\newcommand{\andtxt}{\text{ and }}
\newcommand{\Ais}{\ensuremath{\overset{A}{=}}}
\newcommand{\Cis}{\ensuremath{\overset{C}{=}}}
\newcommand{\closure}[1]{\ensuremath{\overline{#1}}}
\newcommand{\encircle}[1]{%
  \tikz[baseline=(char.base)]{%
    \node[shape=circle, draw, inner sep=1pt] (char) {#1};%
  }%
}
\newcommand{\mathif}{\text{if }}
\newcommand{\seq}[2][n]{\ensuremath{(#2 _{#1})}}
\newcommand{\zz}[1]{$\mathbb{Z}/ #1 \mathbb{Z}$}

\maketitle

\tableofcontents

\pagebreak
\section{Problem 1}
    Let $S$ be the set of all \textit{binary sequences}, that is $S = \{(x_1, x_2, x_3, \dots ) : x_j \in \{0, 1\} \text{ for all } j\}$. 
    \begin{enumerate}[label=(\alph*)]
        \item 
        Using Cantor's diagonalization argument, show that $S$ is uncountable.

        \item 
        Let $T \subset S$ be the set of all binary sequences with \textit{finite support}, that is let \\$T = \{(x_1, x_2, x_3, \dots) \in S : \text{only finitely many of the $x_j$ are nonzero}\}.$ Show that $T$ is countable.
        
        \textit{Hint:} Let $S_n = \{(x_1, x_2, x_3, \dots ) \in S : x_j = 0$ \text{ for all } $j \geq n + 1\}$.
        
    \end{enumerate}

    \subsection{Solution (a)}
    \textbf{Alvin's Proposition: } Define $r: \{0, 1\} \to \{0, 1\}$ where $r(n) = \begin{cases} 0 &\text{ if } n = 1\\1 &\text{ if } n = 0 \end{cases}$.
    For any $x \in \{0, 1\}$, $x \neq r(x)$.
    \begin{proof}[Proof by Cases]
        We consider the cases of $x = 0$ and $x = 1$.

        \textbf{\underline{Case 1: x = 0.}}
        If $x = 0$, then $r(x) = r(0) = 1$.
        Therefore, $x \neq r(x)$.

        \textbf{\underline{Case 2: x = 1.}}
        If $x = 1$, then $r(x) = r(1) = 0$.
        Therefore, $x \neq r(x)$.

        Ergo, for any $x \in \{0, 1\}$, $x \neq r(x)$.
    \end{proof}
    This proof is done through Cantor's Diagonalization Argument, which is a proof by contradiction.
    \begin{proof}[Proof by Contradiction]
        Suppose for contradiction that $S$ is countable.
        That means that we can enumerate the members of $S$.
        We can tabulate these.
        \begin{center}
            \begin{tabular}{ c  c  c  c  c  c  c }
                $S_1 =$ &($x_{11}$,  &$x_{12}$,   &$x_{13}$,   &$x_{14}$,   &$x_{15}$,   &\dots )\\
                $S_2 =$ &($x_{21}$,  &$x_{22}$,   &$x_{23}$,   &$x_{24}$,   &$x_{25}$,   &\dots )\\
                $S_3 =$ &($x_{31}$,  &$x_{32}$,   &$x_{33}$,   &$x_{34}$,   &$x_{35}$,   &\dots )\\
                $S_4 =$ &($x_{41}$,  &$x_{42}$,   &$x_{43}$,   &$x_{44}$,   &$x_{45}$,   &\dots )\\
                $S_5 =$ &($x_{51}$,  &$x_{52}$,   &$x_{53}$,   &$x_{54}$,   &$x_{55}$,   &\dots )\\
                $\vdots$&$\vdots$   &$\vdots$   &$\vdots$   &$\vdots$   &$\vdots$   &$\ddots$
            \end{tabular}
        \end{center}
        
        Note the binary sequence elements which can be written in our enumeration as $x_{jj}$ for some natural number $j$.
        \begin{center}
            \begin{tabular}{ c  c  c  c  c  c  c }
                $S_1 =$ &(\encircle{$x_{11}$},  &$x_{12}$,   &$x_{13}$,   &$x_{14}$,   &$x_{15}$,   &\dots )\\
                $S_2 =$ &($x_{21}$,  &\encircle{$x_{22}$},   &$x_{23}$,   &$x_{24}$,   &$x_{25}$,   &\dots )\\
                $S_3 =$ &($x_{31}$,  &$x_{32}$,   &\encircle{$x_{33}$},   &$x_{34}$,   &$x_{35}$,   &\dots )\\
                $S_4 =$ &($x_{41}$,  &$x_{42}$,   &$x_{43}$,   &\encircle{$x_{44}$},   &$x_{45}$,   &\dots )\\
                $S_5 =$ &($x_{51}$,  &$x_{52}$,   &$x_{53}$,   &$x_{54}$,   &\encircle{$x_{55}$},   &\dots )\\
                $\vdots$&$\vdots$   &$\vdots$   &$\vdots$   &$\vdots$   &$\vdots$   &$\ddots$
            \end{tabular}
        \end{center}

        Let's create a new binary sequence $S_\alpha$.
        Define $r: \{0, 1\} \to \{0, 1\}$ where $r(n) = \begin{cases} 0 &\text{ if } n = 1\\1 &\text{ if } n = 0 \end{cases}$.
        By Alvin's Proposition, for any $x \in \{0, 1\}$, $x \neq r(x)$.
        Take each element written in our enumeration as $x_{jj}$ for some natural number $j$. 
        Define each member $x_{\alpha j}$ of $S$ as $r(x_{jj})$.
        \begin{equation}
            S_\alpha    =   (x_{\alpha 1}, x_{\alpha 2}, x_{\alpha 3}, x_{\alpha 4}, x_{\alpha 5}, \dots)
                        =   (r(x_{11}), r(x_{22}), r(x_{33}), r(x_{44}), r(x_{55}), \dots)
        \end{equation}

        Since for any $x \in \{0, 1\}$, $x \neq r(x)$, that means that for any natural number $k$, $x_{kk} \neq x_{\alpha k}$.
        Therefore, for any binary sequence $S_k$, there exists some element $x_{kk}$ that is different from its respective element $x_{\alpha k}$ of $S_\alpha$.
        That means that $S_\alpha$ is a member of $S$ that differs from each $S_k$ by at last one digit.
        That means that $S_{\alpha}$ cannot occur in our enumeration.
        This contradicts that $S$ is the set of all binary sequences.
        Therefore, our initial assumption is false and $S$ is uncountable. 
    \end{proof}

    \subsection{Solution (b)}
    \begin{proof}
        Recall that the rational numbers are countable.
        This means that if we can inject the binary sequences into the rational numbers, we can prove that they are countable.
        
        Since only finitely many of $x_j$ are nonzero, there is a point $n > 0$ beyond which for all $j > n$, $x_j = 0$.
        Take any element $t$ in $T$ such that $t = (x_1, x_2, \dots)$.
        We can define a number $0 \leq N < 1$, turning every $x_j$ into a digit.
        Since for all $j > n$, $x_j = 0$, there is a terminal point in $N$ beyond which every subsequent digit would be $0$.
        \begin{equation}
            N = 0.x_1 x_2 x_3 \dots x_{n - 1} x_n 0000\dots
        \end{equation}

        More formally, we can define the function from $T$ to $\mathbb{Q}$ as follows.
        \begin{gather}
            f: T \to \mathbb{Q}\\
            f((x_1, x_2, \dots)) = 0.x_1 x_2 x_3 \dots
        \end{gather}
        Since $f$ has a trailing series of zeroes beyond a certain point, the decimal does terminate.
        Note that the number that comes out of $f$ always uses the terminating point form of a number and never the repeating-$9$ expansion.

        We can multiply both sides of $N$ by $10^n$.
        This allows us to shift the decimal point to just after $x_n$.
        \begin{align}
            N \times 10^n   &=  10^n \times 0.x_1 x_2 x_3 \dots x_{n - 1} x_n 0000\dots\\
                &=  x_1 x_2 x_3 \dots x_{n - 1} x_n .0000\dots
        \end{align}
        
        Since it has nothing but $0$ past the decimal point, that means that $N \times 10^n$ is a natural number.
        Furthermore, since $n > 0$, $10^n$ is a natural number.
        We can divide both sides by $10^n$.
        \begin{gather}
            \frac{N \times 10^n}{10^n} = \frac{x_1 x_2 x_3 \dots x_{n - 1} x_n .0000\dots}{10^n}\\
            N   =   \frac{x_1 x_2 x_3 \dots x_{n - 1} x_n}{10^n}
        \end{gather}

        Therefore, $N$ is an integer over a natural number, meaning that it is a rational number.
        In order to establish that this is an injection, we must establish that for $s$ and $t$ in $T$ where $s \neq t$, then $N_s \neq N_t$.
        If $s \neq t$, then there is a smallest number $k$ where $s_k \neq t_k$.
        That leads to digit $k$ of $N_s$ not being equal to digit $k$ of $N_t$.
        This means that $N_s \neq N_t$.

        This establishes an injection from the binary sequences into the rational numbers.
        Therefore, because the rational numbers are countable, $T$ is countable.
    \end{proof}

    \pagebreak

\section{Problem 2}
    Let $(X, d)$ be a metric space. For a subset $E \subset X$, define the \textit{closure} of $E$, denoted $\overline{E}$, to be $E \cup E'$.
    \begin{enumerate}[label=(\alph*)]
        \item 
        Show that $\overline{E}$ is closed, and that $\overline{E} = E$ if and only if $E$ is closed.

        \item 
        Show that $\overline{E}$ is the \textit{smallest} closed set containing $E$ in the following sense: if $F$ is closed and $E \subset F$, then $\overline{E} \subset F$.

        \item 
        Show that $\overline{E}$ is the intersection of all closed subsets of $X$ containing $E$.

        \item 
        Show that for any $E, F \subset X$, we have $\overline{E \cup F} = \overline{E} \cup \overline{F}$ (this is much easier if you use part (b)).

        \item 
        Define the \textit{closed ball of radius $r$ centered at $x$} to be $D_r(x) = \{y \in X : d(y, x) \leq r\}$. Show that $D_r(x)$ is closed and that $\overline{B_r(x)} \subset D_r(x)$ (using part (b) is again much simpler here).
        
    \end{enumerate}

    \subsection{Solution (a)}
        First we prove that $\overline{E}$ is closed.
        \begin{proof}
            We do this by proving that the complement of \closure{E} is open.
            Take any $g$ in $\closure{E}^c$.
            Since $\closure{E}$ is equal to the union of $E$ and the set of all its limit points, $g$ is neither a limit point of $E$ nor in $E$.
            Not all $r > 0$ create an open ball $B_r(g)$ that has an intersection with \closure{E}.
            That means that there exists some $r > 0$ such that the open ball $B_r(g)$ is a subset of $\closure{E}^c$.
            This means that $\closure{E}^c$ is open.
            Because a set is open if and only if its complement is closed, \closure{E} is closed.
        \end{proof}
        Next, we prove that $\closure{E} = E$ if $E$ is closed.
        \begin{proof}
            If $E$ is closed, that means that it contains all its limit points.
            Recall that $\closure{E} = E \cup E'$, with $E'$ referring to the set of all limit points of $E$.
            This means that $E'$ is a subset of $E$.
            Therefore, $E \cup E'$ is equal to $E$.
            Ergo, transistively, $\closure{E} = E$.
        \end{proof}
        Third, we prove that $E$ is closed if $\closure{E} = E$.
        \begin{proof}
            If $\closure{E} = E$, then $E$ is equal to the union of $E$ and the set of all its limit points $E'$.

            Take any point $p$ in $E'$.
            Since $p$ is in $E'$, that means $p$ is in the union of $E$ and $E'$.
            Since $E$ is equal to the union of $E$ and $E'$, that means that $p$ is in $E$.
            This means that if a point is a limit point of $E$, then it is in $E$.
            That means that $E$ contains all of its limit points.
            Ergo, $E$ is closed.
        \end{proof}
        Last, unification.
        \begin{proof}
            We know that $\closure{E} = E$ if $E$ is closed.
            We know that $E$ is closed if $\closure{E} = E$.
            Ergo, $\closure{E} = E$ if and only if $E$ is closed.
        \end{proof}
    
    \subsection{Solution (b)}
        \begin{proof}
            We will prove that if $F$ is closed and $E$ is a subset of $F$, then \closure{E} is a subset of $F$.

            Observe the limit points of $E$ (denoted $E'$), whether they exist or not.
            Take any point $p$ in $E'$.
            Since $p$ is a limit point of $E$, that means that for every $r > 0$, there exists an element $q$ in $E$ that is not equal to $p$ such that the open ball $B_r(p)$ contains $q$.
            Since $E$ is a subset of $F$, this means that if $q$ is in $E$, then $q$ is also in $F$.
            Therefore, for every $r > 0$, there exists an element $q$ in $F$ that is not equal to $p$ such that the open ball $B_r(p)$ contains $q$.
            This means that $p$ is a limit point of $F$, and therefore that every limit point of $E$ is a limit point of $F$.
            Since $F$ is closed, that means it contains all its limit points.
            Therefore $E'$ is a subset of $F$.
            Since $E$ and $E'$ are both subsets of $F$, that means their union is also a subset of $F$.
            Since \closure{E} is the union of $E$ and $E'$, that means \closure{E} is a subset of $F$.
        \end{proof}

    \subsection{Solution (c)}
        \begin{proof}
            In part (b), we proved that if $F$ is closed and $E$ is a subset of $F$, then \closure{E} is a subset of $F$.
            We also know that the intersection of any number of closed sets is also a closed set.
            Supppose that we give the set of all closed subsets of $X$ containing $E$ the name $G$.

            The only thing we know all the sets in $G$ have in common is that they contain $E$ as a subset. 
            Therefore, we can say that the intersection of all elements of $G$ contains $E$ as a subset.
            As proven in part (b), if a closed set has $E$ as a subset, then \closure{E} is a subset of the closed set.
            Therefore, \closure{E} is a subset of the intersection of all elements of $G$.

            Since $G$ is the collection of all closed subsets of $X$ and we have established that \closure{E} is both closed and a subset of $X$, we know that \closure{E} is an element of $G$.
            By that, we know that if an element $p$ of $X$ is not in \closure{E}, it is not in the intersection of all elements of $G$.
            Therefore, the intersection of all elements of $G$ is a subset of \closure{E}.

            Since the intersection of all elements of $G$ is both a subset and superset of \closure{E}, that means the intersection of all elements of $G$ is equal to \closure{E}.
        \end{proof}

    \subsection{Solution (d)}
        \textbf{Baron's Proposition:} $A \subset C$ and $B \subset C$ if and only if $A \cup B \subset C$.
        \begin{proof}
            First the forward direction.
            Suppose that $A \subset C$ and $B \subset C$.
            Any element in $A \cup B$ is either in $A$ or in $B$.
            If it is in $A$, since $A \subset C$, it is also in $C$.
            If it is in $B$, since $B \subset C$, it is also in $C$.
            Either way, it is in $C$, so every element in $A \cup B$ is in $C$.
            Therefore, $A \cup B$ is a subset of $C$.

            Next the backwards direction.
            Suppose that $A \cup B \subset C$.
            Every element in $A$ is in $A \cup B$.
            Since $A \cup B \subset C$, every element in $A$ is in $C$.
            Similarly, every element in $B$ is in $A \cup B$.
            Since $A \cup B \subset C$, every element in $B$ is in $C$.
            Therefore, $A \subset C$ and $B \subset C$.

            Ergo, $A \subset C$ and $B \subset C$ if and only if $A \cup B \subset C$.
        \end{proof}

        First, we demonstrate that \closure{E \cup F} is a subset of $\closure{E} \cup \closure{F}$.
        \begin{proof}
            Since the union of finite closed sets is also a closed set, $\closure{E} \cup \closure{F}$ is a closed set.
            Since \closure{E} is equal to the union of $E$ and $E'$, then $E$ is a subset of \closure{E}.
            This means that $E$ is a subset of $\closure{E} \cup \closure{F}$.
            By a similar argument, $F$ is a subset of $\closure{E} \cup \closure{F}$.
            Since $E$ and $F$ are both subsets of $\closure{E} \cup \closure{F}$, Baron's Proposition tells us that $E \cup F$ is a subset of $\closure{E} \cup \closure{F}$.
            Since $\closure{E} \cup \closure{F}$ is a closed set, according to part (b), \closure{E \cup F} is also a subset of $\closure{E} \cup \closure{F}$.
        \end{proof}

        Next, we demonstrate that $\closure{E} \cup \closure{F}$ is a subset of \closure{E \cup F}.
        \begin{proof}
            First, we can write out \closure{E \cup F} in long form.
            \begin{equation}
                \closure{E \cup F} = (E \cup F) \cup (E \cup F)'
            \end{equation}

            Since it is written as a union including $E$, we can say that $E$ is a subset of \closure{E \cup F}.
            Since \closure{E \cup F} is closed, by part (b), \closure{E} is a subset of \closure{E \cup F}.
            By a similar argument, \closure{F} is a subset of \closure{E \cup F}.
            
            By Baron's Proposition, that means that $\closure{E} \cup \closure{F}$ is a subset of \closure{E \cup F}.
        \end{proof}

        Last, we combine the two subset statements.
        \begin{proof}
            We know that $\closure{E} \cup \closure{F}$ is a subset of \closure{E \cup F} and that \closure{E \cup F} is a subset of $\closure{E} \cup \closure{F}$.
            Since each is a subset the other, that means that $\closure{E \cup F} = \closure{E} \cup \closure{F}$.
        \end{proof}

    \subsection{Solution (e)}
        First we prove that $D_r(x)$ is closed.
        \begin{proof}
            For $D_r(x)$ to be closed, its complement would have to be open.
            We can define its complement.
            \begin{equation}
                D_r(x)^c = \{ y \in X : d(y,x) > r \}
            \end{equation}

            Take any point $p$ in $D_r(x)^c$.
            We know that $d(p,x) > r$.
            That means that $d(p,x) - r > 0$.
            Suppose that we create an open ball of radius $s = d(p,x) - r$.
            \begin{equation}
                B_s(p) = \{ y \in X : d(y,p) < d(p,x) - r\}
            \end{equation}

            If we can prove that $d(y,x) > r$ for any $y \in B_s(p)$, we can prove that the open ball is a subset of $D_r(x)^c$.
            By the triangle inequality, with $x$ as our third point, $d(x,p) \leq d(y,x) + d(y,p)$.
            This means that $d(y,x) \geq d(x,p) - d(y,p)$.
            Since $d(p,x) > r$, we know $d(y,x) > d(x,p) - d(x,p) + r = r$.

            Therefore, every point in $B_s(p)$ is in $D_r(x)^c$.
            This means that $D_r(x)^c$ is open.
            Ergo, $D_r(x)$ is closed.
        \end{proof}

        We can now prove that $\closure{B_r(x)}$ is a subset of $D_r(x)$.
        \begin{proof}
            First, we prove that $B_r(x)$ is a subset of $D_r(x)$.
            Take any point $y$ in $B_r(x)$.
            We can say that $d(x,y) < r$.
            That means that $d(x,y) \leq r$.
            Therefore, $y$ is in $D_r(x)$.
            This means that $B_r(x)$ is a subset of $D_r(x)$.

            Part (b) tells us that if $F$ is closed and $E$ is a subset of $F$, then \closure{E} is a subset of $F$.
            $D_r(x)$ is closed.
            $B_r(x)$ is a subset of $D_r(x)$.
            Therefore, \closure{B_r(x)} is a subset of $D_r(x)$.
        \end{proof}
    \pagebreak

\section{Problem 3}
    Let $S$ be the set of all four letter strings of letters (which we will call \textit{words}). Thus $S = \{\text{abcd, eeee, math, wyzx, ...}\}$. Define a function $d$ which takes in two words and outputs the number of positions at which they differ. So for example, $d(\text{math, cats}) = 2$ and $d(\text{abcd}, \text{dcba}) = 4$.

    \begin{enumerate}[label=(\alph*)]
        \item
        Show that $(S, d)$ is a metric space (this is called the \textit{Hamming metric}).

        \item 
        In $(S, d)$, consider the sets
        \begin{align*}
            B_2(\text{aaaa}) &= \{s \in S : d(s, \text{aaaa}) < 2)\} \quad \text{ and }\\[5pt] D_2(\text{aaaa}) &= \{s \in S : d(s, \text{aaaa}) \leq 2)\}.
        \end{align*}
        Is it true that the closure of $B_2(\text{aaaa})$ equals $D_2(\text{aaaa})$? Using the results from Problem 2 will make things easier.
        
    \end{enumerate}

    \subsection{Solution (a)}
        A metric space consists of two parts: a set and a distance function.
        The metric space is $S$.
        The function is $d$.
        We need to prove that the distance function satisfies three conditions for some $p$ and $q$ in $S$:
        \begin{enumerate}[label=(\alph*)]
            \item   $d(p,q) > 0$ if $p \neq q$; $d(p,q) = 0$ if $p = q$
            \item   $d(p,q) = d(q,p)$
            \item   $d(p,q) \leq d(p,r) + d(r,q)$, for any $r$ in $S$
        \end{enumerate}

        Start with the first.
        \begin{proof}
            We can view words as a finite sequences of letters ($p = \{p_1, p_2, p_3, p_4\}$ and $q = \{q_1, q_2, q_3, q_4\}$).
            We can view the distance function as an alteration of the discrete function ($f$).
            \begin{gather}
                f: X \to \{0, 1\}\\
                f(a,b) = \begin{cases}
                    0   &\text{if } x = y\\
                    1   &\text{if } x \neq y
                \end{cases}\\
                d(p,q) = \sum_{i = 1}^{4} f(p_i, q_i)
            \end{gather} 

            Suppose that $p \neq q$.
            Therefore, they differ in at least one position.
            Since they differ in at least one position, the distance function would be greater than or equal to $1$.
            Therefore, it would be greater than zero.

            Suppose $p = q$.
            There would be no positions in which the two differ.
            \begin{equation}
                d(p,q) = 0 + 0 + 0 + 0 = 0
            \end{equation}
            
            Therefore, the first is satisfied.
        \end{proof}

        Now the second ($d(p,q) = d(q,p)$).
        \begin{proof}
            Recall our equation for the distance function in $(S,d)$.
            \begin{gather}
                f: X \to \{0, 1\}\\
                f(a,b) = \begin{cases}
                    0   &\text{if } x = y\\
                    1   &\text{if } x \neq y
                \end{cases}\\
                d(p,q) = \sum_{i = 1}^{4} f(p_i, q_i)
            \end{gather} 
            
            $f$ is commutative (since $x = y$ means $y = x$ and $x \neq y$ means $y \neq x$).
            Therefore, we can say that $d(p,q) = \sum_{i = 1}^{4} f(p_i, q_i) = \sum_{i = 1}^{4} f(q_i, p_i) = d(q,p)$.
            Ergo, $d(p,q) = d(q,p)$.
        \end{proof}

        Now the third (the triangle inequality).
        \begin{proof}
            The triangle inequality says that for any $r$ in $S$, $d(p,q) \leq d(p,r) + d(r,q)$.
            We can turn the two distance functions into the sums that we have described them as.
            Since the functions are finite, we know we can use the associative property to combine them into one.
            \begin{equation}
                d(p,r) + d(r,q) = \sum_{i = 1}^{4} f(p_i, r_i) + \sum_{i = 1}^{4} f(r_i, q_i)
                =   \sum_{i = 1}^{4} f(p_i, r_i) + f(r_i, q_i)
            \end{equation}

            For $r = \{r_1, r_2, r_3, r_4\}$, there are five possibilities for each letter in the sequence: 
            \begin{enumerate}
                \item   $r_i = p_i$ and $p_i \neq q_i$
                \item   $r_i = q_i$ and $p_i \neq q_i$
                \item   $p_i = q_i$ and $r_i \neq p_i$
                \item   $r_i = p_i = q_i$
                \item   $r_i \neq p_i$, $p_i \neq q_i$, and $r_i \neq q_i$
            \end{enumerate}
            We will prove that $f(p_i,q_i) \leq f(p_i, r_i) + f(r_i, q_i)$ in all four cases.

            \textbf{Case 1: $r_i = p_i$ and $p_i \neq q_i$.}
            Transistively, we can say that $r_i \neq q_i$.
            $p_i \neq q_i$, so $f(p_i,q_i) = 1$.
            $p_i = r_i$, so $f(p_i,r_i) = 0$.
            $r_i \neq q_i$, so $f(p_i,q_i) = 1$.
            Therefore,  $f(p_i, r_i) + f(r_i, q_i) = 1$, so $f(p_i,q_i) \leq f(p_i, r_i) + f(r_i, q_i)$.

            \textbf{Case 2: $r_i = q_i$ and $p_i \neq q_i$.}
            Transistively, we can say that $r_i \neq p_i$.
            $p_i \neq q_i$, so $f(p_i,q_i) = 1$.
            $q_i = r_i$, so $f(q_i,r_i) = 0$.
            $p_i \neq r_i$, so $f(p_i,r_i) = 1$.
            Therefore,  $f(p_i, r_i) + f(r_i, q_i) = 1$, so $f(p_i,q_i) \leq f(p_i, r_i) + f(r_i, q_i)$.

            \textbf{Case 3: $p_i = q_i$ and $r_i \neq p_i$.}
            Transistively, we can say that $r_i \neq q_i$.
            $p_i = q_i$, so $f(p_i,q_i) = 0$.
            $q_i \neq r_i$, so $f(p_i,r_i) = 1$.
            $p_i \neq r_i$, so $f(p_i,r_i) = 1$.
            Therefore,  $f(p_i, r_i) + f(r_i, q_i) = 2$, so $f(p_i,q_i) \leq f(p_i, r_i) + f(r_i, q_i)$.

            \textbf{Case 4: $r_i = p_i = q_i$.}
            $p_i = q_i$, so $f(p_i,q_i) = 0$.
            $q_i = r_i$, so $f(p_i,r_i) = 0$.
            $p_i = r_i$, so $f(p_i,r_i) = 0$.
            Therefore,  $f(p_i, r_i) + f(r_i, q_i) = 0$, so $f(p_i,q_i) \leq f(p_i, r_i) + f(r_i, q_i)$.

            \textbf{Case 5: $r_i \neq p_i$, $p_i \neq q_i$, and $r_i \neq q_i$.}
            $p_i \neq q_i$, so $f(p_i,q_i) = 1$.
            $q_i \neq r_i$, so $f(p_i,r_i) = 1$.
            $p_i \neq r_i$, so $f(p_i,r_i) = 1$.
            Therefore,  $f(p_i, r_i) + f(r_i, q_i) = 2$, so $f(p_i,q_i) \leq f(p_i, r_i) + f(r_i, q_i)$.

            Since in every case, $f(p_i,q_i) \leq f(p_i, r_i) + f(r_i, q_i)$, we can say the same about their sum.
            \begin{equation}
                \sum_{i = 1}^{4} f(p_i,q_i) \leq \sum_{i = 1}^{4} f(p_i, r_i) + f(r_i, q_i)
                    =   \sum_{i = 1}^{4} f(p_i, r_i) + \sum_{i = 1}^{4} f(r_i, q_i)
            \end{equation}

            We can replace the sums with the distance functions, so $d(p,q) \leq d(p,r) + d(r,q)$.
        \end{proof}

        Lastly, we consolidate this and formally prove that $(S,d)$ is a metric space.
        \begin{proof}
            A metric space consists of two parts: a set and a distance function.
            The set is $S$.
            The function is $d$.
            We need to prove that the distance function satisfies three conditions for some $p$ and $q$ in $S$:
            \begin{enumerate}[label=(\alph*)]
                \item   $d(p,q) > 0$ if $p \neq q$; $d(p,q) = 0$ if $p = q$
                \item   $d(p,q) = d(q,p)$
                \item   $d(p,q) \leq d(p,r) + d(r,q)$, for any $r$ in $S$
            \end{enumerate}

            We have earlier established that the distance function satisfies all three conditions.
            Ergo, $(S,d)$ is a metric space.
        \end{proof}

    \subsection{Solution (b)}
        I will use Problem 2, part (e).
        It is false that the closure of $B_2(\text{aaaa})$ is equal to $D_2(\text{aaaa})$.
        \begin{proof}
            $D_2(\text{aaaa})$ fits the definition of a closed ball in $S$, which is found in the instructions for Problem 2. 

            Take any $b$ in $B_2(\text{aaaa})$.
            This means that $d(b, \text{aaaa}) < 2$.
            Therefore, $d(b, \text{aaaa}) \leq 2$.
            This means that $b$ is in $D_2(\text{aaaa})$.
            Therefore, $B_2(\text{aaaa})$ is a subset of $D_2(\text{aaaa})$.

            The issue comes when we consider the cases where $d(b, \text{aaaa}) = 2$.
            We can prove that the cases where $d(b, \text{aaaa}) = 2$ are not limit points of $B_2(\text{aaaa})$.
            In part (a), we showed that if $p \neq q$, then $d(p,q) \geq 1$.
            We also showed that $d(p,q) = 0$ if $p = q$.
            Take any point $b$ for which $d(b, \text{aaaa}) = 2$.
            Suppose some ball of radius $r > 0$ around $b$ such that $r < 1$.
            Since $d(p,q) \geq 1$ if $p \neq q$, then $p = q$ if $d(p,q) < 1$.
            If $p = q$, then $d(p,q) = 0$.
            Therefore, if $d(p,q) < 1$, then $d(p,q) = 0$.
            If $d(p,q) < r$, then $d(p,q) < 1$.
            Therefore, if $d(p,q) < r$, then $d(p,q) = 0$ and $p = q$.
            This means that if $r < 1$, then the only point $p$ where $d(p,b) < r$ is when $p = b$.
            Therefore, if $r < 1$, then $B_r(b)$ has no point in $B_2(\text{aaaa}) \backslash \{b\}$.
            Therefore, $b$ is not a limit point of $B_2(\text{aaaa})$.
            Therefore, the points $b$ where $d(b, \text{aaaa}) = 2$ are not limit points of $B_2(\text{aaaa})$.

            Since the points $b$ where $d(b, \text{aaaa}) = 2$ are not limit points of $B_2(\text{aaaa})$, they are not in \closure{B_2(\text{aaaa})}.
            However, if $d(b, \text{aaaa}) = 2$, then $d(b, \text{aaaa}) \leq 2$.
            Therefore, the points are in $D_2(\text{aaaa})$ but not \closure{B_2(\text{aaaa})}.
            Ergo, while it is a subset, $\closure{B_2(\text{aaaa})} \neq D_2(\text{aaaa})$.
        \end{proof}
    \pagebreak

\section{Problem 4}
    Let $E = \mathbb{Q} \cap [0, 1]$, viewed as a subset of $(\mathbb{R}, d)$, where $d$ is the usual metric.
    \begin{enumerate}[label=(\alph*)]
        \item 
        Find an open cover of $E$ that has no finite subcover, and conclude that $E$ is not compact. (\textit{Hint:} pick an irrational number $r \in [0, 1]$ to help you.)

        \item 
        Find a set $A \subset E$ (other than $\varnothing$ or $E$) that is clopen w.r.t. the subspace metric on $(E, d)$, and conclude that $E$ is not connected. You may again find picking an irrational $r \in [0, 1]$ helpful.

        \item 
        With respect to $(\mathbb{R}, d)$, is the set $A$ you chose in part (b) open, closed, both, or neither? Explain.
        
    \end{enumerate}
    
    \subsection{Solution (a)}
        \textbf{Chloe's Proposition:}
            Rational numbers are closed under addition.
        \begin{proof}
            Take two rational numbers $\frac{m}{n}$ and $\frac{p}{q}$ where $m$ and $p$ are integers and $n$ and $q$ are natural numbers (not including zero).
            Adding them together, we can cross multiply and combine the two.
            \begin{equation}
                \frac{m}{n} + \frac{p}{q}   =   \frac{mq}{nq} + \frac{np}{nq}
                    =   \frac{mq + np}{nq}
            \end{equation}

            Since integers are closed under multiplication and addition, and the natural numbers are a subset of the integers, $mq + np$ is an integer, and $nq$ is a natural number.
            Therefore, $\frac{mq + np}{nq}$ is a rational number.
        \end{proof}
        
        \textbf{Dylan's Proposition:}
            The sum of a rational number and an irrational number is an irrational number.
        \begin{proof}[Proof by Contradiction]
            Suppose for contradiction that for irrational number $r$ and rational number $p$, their sum is a rational number $q$.
            \begin{equation}
                r + p = q
            \end{equation}

            Subtract $q$ from both sides.
            \begin{equation}
                r = q - p
            \end{equation}

            Since rational numbers are closed under addition (see Chloe's Proposition) and both $p$, $q$, and $-p$ are integers, $q - p$ is an integer and therefore $r$ is an integer.
            Since the integers are the subsets of the rational numbers, that means $r$ is a rational number.
            This contradicts the fact that $r$ is irrational.
            Therefore our initial assumption is false and the sum of a rational number and an irrational number is an irrational number.
        \end{proof}
        
        \textbf{Eddie's Proposition:}
            In an ordered field, if $a < b$, then $-a > -b$.
        \begin{proof}
            Suppose $a < b$.
            Add $-a - b$ to both sides to get $a - a - b < b - a - b$.
            By the additive inverse of both $a$ and $b$, $-b < -a$.
        \end{proof}
        
        \textbf{Frank's Proposition:}
            $\abs{a} < b$ if and only if $-b < a < b$.
        \begin{proof}
            Suppose $\abs{a} < b$.
            This divides us into two cases: $a \geq 0$ and $a < 0$.
            
            \textbf{Case 1: $a \geq 0$.}
                If $a \geq 0$, there are two results.
                First, through Eddie's Proposition and because $a \geq 0$, $-a \leq 0$.
                Second, by the definition of the absolute value, $\abs{a} = a$, so $a < b$.
                Also by Eddie's Proposition, $-b < -a$.
                Transistively through all of these, $-b < -a \leq 0 \leq a < b$.
                Shortening this, we get $-b < a < b$.

            \textbf{Case 2: $a < 0$.}
                There are two results from $a < 0$.
                First, through Eddie's Proposition and because $a < 0$, $-a > 0$.
                Second, by the definition of the absolute value, $\abs{a} = -a$, so $-a < b$.
                Also by Eddie's Proposition, $-b < a$.
                Transistively again, $-b < a < 0 < -a < b$.
                Therefore, $-b < a < b$.
        \end{proof}

        Now we prove the open cover problem.
        \begin{proof}
            Take an irrational number $r$ in the set $[0,1] \backslash \mathbb{Q}$. 

            Let us create two sets of open intervals, and call them $L$ and $U$, and call their full union $C$.
            \begin{gather}
                L = \left\{ \left( r - \frac{1}{n}, r - \frac{1}{n + 1} \right) \Big| n \in \mathbb{N} \right\}\\
                U = \left\{ \left( r + \frac{1}{n + 1}, r + \frac{1}{n} \right) \Big| n \in \mathbb{N} \right\}\\
                C = \bigcup_{n \in \mathbb{N}} L_n \cup \bigcup_{n \in \mathbb{N}} U_n
            \end{gather}

            Since they are made up of unions of open sets, the union of all elements of $L$ and $U$ are also open sets.

            We can prove that $C$ is a cover of $[0,1] \cap \mathbb{Q}$.
            Take any element of $[0,1] \cap \mathbb{Q}$ and call it $p$.
            Since $r$ is irrational, it is not in $[0,1] \cap \mathbb{Q}$, so either $p < r$ or $p > r$.
            We can prove that $p$ is in some element in $L$ or $U$ by cases.

            \textbf{Case 1: $p < r$.}
                Since $p \in [0,1]$, we know that $0 \leq p$.
                Since $p < r$, $r - p > 0$.
                By the Archimedean Principle, there is some natural number $n$ such that $\frac{1}{n} > r - p \geq \frac{1}{n + 1}$.
                Since $r$ is irrational and $p$ is rational, according to Dylan's Proposition, $r - p$ is irrational.
                $r - p$ is irrational but $\frac{1}{n + 1}$ is rational, so we can safely say $r - p \neq \frac{1}{n + 1}$, and thereby $\frac{1}{n} > r - p > \frac{1}{n + 1}$.
                Multiply everything by $-1$ (and flip all inequality signs, which is done due to Eddie's Proposition) to find that $-\frac{1}{n} < p - r < -\frac{1}{n + 1}$.
                Add $r$ to all sides to find that $r - \frac{1}{n} < p < r - \frac{1}{n + 1}$.
                Therefore, $p$ is in $\left( r - \frac{1}{n}, r - \frac{1}{n + 1} \right)$ for some natural number $n$.

            \textbf{Case 2: $p > r$.}
                Since $p \in [0,1]$, we know that $1 \geq p$.
                Since $p > r$, $p - r > 0$.
                By the Archimedean Principle, there is some natural number $n$ such that $\frac{1}{n} > p - r \geq \frac{1}{n + 1}$.
                Since $r$ is irrational and $p$ is rational, according to Dylan's Proposition, $p - r$ is irrational.
                $p - r$ is irrational but $\frac{1}{n + 1}$ is rational, so we can safely say $r - p \neq \frac{1}{n + 1}$, and thereby $\frac{1}{n} > p - r > \frac{1}{n + 1}$.
                Add $r$ to all sides to find that $r + \frac{1}{n} > p > r + \frac{1}{n + 1}$.
                Therefore, $p$ is in $\left( r + \frac{1}{n + 1}, r + \frac{1}{n} \right)$ for some natural number $n$.

            That means that $C$ is a cover of $[0,1] \cap \mathbb{Q}$.
            Since all sets in $L$ and $U$ are open, it is an open cover of $[0,1] \cap \mathbb{Q}$.
            If we take a finite subset of $C$, that means that there will be some largest index in the creation of both $L$ and $U$, which we can call $N_L$ and $N_U$ respectively.
            Therefore, there will be an interval of $\left( r - \frac{1}{N_L + 1}, r + \frac{1}{N_U + 1} \right)$ that is not in the open subset.
            Therefore, the subset is not a cover of $[0,1] \cap \mathbb{Q}$.
            Therefore, $[0,1] \cap \mathbb{Q}$ has a cover with no finite subcover and is not compact.
        \end{proof}

    \subsection{Solution (b)}
        \begin{proof}[Proof by Counterexample]
            Take the case of, for the irrational number $r$ where $0 < r < 1$, $A = [0,r)$.
            To prove this is clopen, we prove it is both open and closed.
            
            \textbf{Case 1: $A$ is open.}
                Take any point $p$ in $A$.
                Create an open ball $B_s(p)$ in $E$ around $p$ of radius $s = r - p$, defined as $B_s(p) = \left\{ x \in E : d(x,p) < s \right\}$.
                Since $r > p$, $s = r - p > 0$.
                Since $d$ is the usual metric, this is the equivalent of saying that $\left| x - p \right| < r - p$.
                Add $p$ to both sides to get $\left| x - p \right| + p < r$.
                The triangle inequality tells us that $\left| x - p + p \right| = \abs{x} < r$.
                Since for all $x$ in $E$, $x \geq 0$, then we can say $x = \abs{x}$, so $0 \leq x < r$.
                Therefore, $x$ is in $A$.
                That means $B_s(p)$ is in $A$.
                This means that $A$ is open.

            \textbf{Case 2: $A$ is closed.}
                If $A$ is closed, then $A^c$ would be open.
                $A$ is defined as $[0,r)$, so $A^c = [r,1]$.
                Since $r$ is irrational, $E$ is a subset of the rational numbers, and $A$ is a subset of $E$, $r$ is not in $A$, so $A = (r,1]$.
                Take any point $p$ in $A$.
                Create an open ball $B_s(p)$ in $E$ around $p$ of radius $s = p - r$, defined as $B_s(p) = \left\{ x \in E : d(x,p) < s \right\}$.
                Since $r < p$ and $p < 1$, $s = p - r < 1 - r$.
                Since $d$ is the usual metric, this is the equivalent of saying that $\left| x - p \right| < p - r$.
                By Frank's proposition, $r - p < x - p < p - r$.
                Add $p$ to all sides to find $r < x < 2p - r$, and we can focus on the former two terms.
                Add $r$ to both sides to get $\left| x - p \right| + r < p$.
                The triangle inequality tells us that $\left| x - p + p \right| = \abs{x} < r$.
                Since for all $x$ in $E$, $x \leq 1$, then we can say $r < x \leq 1$.
                Therefore, $x$ is in $A^c$.
                That means $B_s(p)$ is in $A^c$.
                This means that $A^c$ is open.
                This means that $A$ is closed.

            Therefore, $A$ is clopen and $A$ is not equal to either $E$ or $\varnothing$.
            Since a set $E \subset \mathcal{X}$ is connected if the only clopen sets in $E$ are $E$ and $\varnothing$, $E$ is not connected.
        \end{proof}

    \subsection{Solution (c)}
        In $(\mathbb{R}, d)$, $A$ is neither.
        First we prove that it is not closed.
        \begin{proof}[Proof by Counterexample]
            For a set to be closed, it must contain all its limit points.
            $r$ is a limit point of $A$ because for any $\varepsilon > 0$ where $\varepsilon < r$, $r - \varepsilon < r$ and $r - \varepsilon > 0$, so $0 < r - \varepsilon < r$, and the same applies to $\frac{\varepsilon}{2}$.
            Since the rationals are dense in the reals, there is some rational number $q$ such that $r - \varepsilon < q < r$.
            Therefore, for any open ball of radius $\varepsilon > 0$, $q$ is in $A$ and in the open ball.

            Ergo, $\mathbb{Q} \cap [0,r)$ does not contain all its limit points.
            This means that $A$ is not closed.
        \end{proof}

        Next, we prove that it is not open.
        \begin{proof}[Proof by Counterexample]
            If a set is open, then every point in it has an open ball around it also contained in the set.
            Take the case of $0$.
            $0$ is in $A$.
            However, for any open ball of radius $\varepsilon > 0$, $-\frac{\varepsilon}{2}$ is in the open ball.
            However, since $-\frac{\varepsilon}{2} < 0$, $-\frac{\varepsilon}{2}$ is not in $A$.
            Therefore, there is no open ball of radius $\varepsilon$ in $\mathbb{R}$ centered at $0$ that is a subset of $A$.
            Therefore, $A$ is not open.
        \end{proof}
    \pagebreak

\section{Problem 5}
    Let $(X, d)$ be an arbitrary metric space.
    \begin{enumerate}[label=(\alph*)]
        \item 
        Show that a finite union of compact sets in $(X, d)$ is compact. Is this also true for arbitrary unions? Prove or give a counterexample.

        \item 
        Show that a set $K \subset (X, d)$ is compact if and only if every cover of $K$ \textit{consisting of open balls} has a finite subcover. (\textit{Hint:} For the trickier direction, given an open cover $\{U_\alpha\}_{\alpha \in A}$ of $K$, make a new cover that consists of all of the open balls contained in some $U_\alpha$.)

        \item 
        Show that compact $K \subset (X, d)$ have the following property: for all $\epsilon > 0$, there is a finite set $\{x_1, \dots, x_n\} \subset K$ so that each $y \in K$ is within distance $\epsilon$ of some $x_j$.
        
    \end{enumerate}
    
    \subsection{Solution (a)}
        First, the finite union.
        \begin{proof}[Proof by Construction]
            We will be using notations and names for the sake of improved comprehension.
            For an open cover $C$ over a union of all sets in $\{S_k\}_{k=1}^N$, we will be constructing a finite subcover $G$.

            A set is compact if and only if every open cover of it has a finite subcover.
            Therefore, since each $S_k$ is compact, $S_k$ has a finite subcover.

            Call the finite collection of $N$ compact sets we will take the union over $\{S_k\}_{k=1}^N$ and call their union $A$.
            Take any open cover of $A$ called $C$.
            Since $A$ is a union of all $\{S_k\}_{k=1}^N$, for any natural number $k < N$, $S_k$ is a subset of $A$.
            Therefore, $C$ is also an open cover of $S_k$.
            Since $S_k$ is compact, $C$ has a finite subcover of $S_k$.
            Call that finite subcover $F_k$.
            Therefore, for every set $S_k$, there is a finite subcover of $S_k$ taken from $C$ called $F_k$.

            Take the union of $F_k$ for all natural numbers $k$ where $1 \leq k \leq N$ and call that union $G$.
            Since each $F_k$ is finite and $N$ is finite, $\bigcup_{k = 1}^{N} F_k$ is finite a finite set.
            For every $F_k$, since $G$ contains every element of $F_k$, $F_k$ is a subset of $G$.
            Since $F_k \subset G \subset C$ and $F_k$ covers $S_k$, $G$ also covers $S_k$.
            Since $G$ is a finite subcover of $C$ that covers every $S_k$ and $A$ is the union of every $S_k$, $G$ is a finite subcover over $A$.
            This means $G$ is a finite subcover of $C$ over $A$.
            Therefore, for every open cover $C$ over $A$, $A$ has a finite subcover.
            Ergo, $A$ is compact.
        \end{proof}

        For an arbitrary union, we can consider a case of infinite compact sets.
        \begin{proof}[Proof by Counterexample]
            Take $(X,d) = (\mathbb{R},d)$ and the set $S = \left\{ [n,n+1] \Big| n \in \mathbb{N} \right\}$.
            By the Bolzano-Weistrass Theorem, since every element of $S$ is in $\mathbb{R}$ and is closed and bounded, they are all compact sets.
            Take the union $A$ of all the sets in $S$.
            The union therefore becomes $[1,\infty)$.
            However, by the Bolzano-Weistrass theorem, since $[1,\infty)$ is not bounded, it is not compact.
        \end{proof}

    \subsection{Solution (b)}
        First we prove that if every cover of $K \subset (X,d)$ consisting of open balls has a finite subcover, then the set $K$ is compact.
        \begin{proof}[Proof by Construction]
            Take any open cover of $K$ called $\mathcal{U}$.
            Take any open set $U \in \mathcal{U}$.
            By the definition of the open set, for any $x \in U$, there is some open ball of radius $r > 0$ centered around $x$ called $B_r(x)$ such that $x \in B_r(x) \subset U \in \mathcal{U}$. 
            Therefore, every $x$ in $K$ is contained within some $B_r(x)$.
            
            Take every $B_r(x)$ for every $U$ in $\mathcal{U}$ and form a new set, which we can call $\mathcal{B}$.
            Since every element of $K$ is contained inside some open ball $B_r(x)$ in $\mathcal{B}$, and because every cover of $K$ consisting of open balls has a finite subcover, $\mathcal{B}$ has a finite subcover of $K$.
            Call that finite subcover $\mathcal{F}$.

            We have previously established that every $B_r(x)$ is a subset of some $U \in \mathcal{U}$.
            Create a new set of open sets $F$.
            For every $B_r(x)$ in $\mathcal{F}$, pick one and only one $U$ in $\mathcal{U}$ such that $B_r(x)$ is a subset of $U$ and add $U$ to $F$.
            Therefore, $F$ has equal or lower cardinality than $\mathcal{F}$, so it is finite and is a subset of $\mathcal{U}$.

            Since $\mathcal{F}$ is an open cover of $K$ and every $B_r(x)$ in $\mathcal{F}$ is a subset of some open set $U$ in $F$, $F$ is also an open cover of $K$.
            Since $F$ is an open cover of $K$, is finite, and is a subset of $\mathcal{U}$, it is a finite open subcover of $\mathcal{U}$ over $K$.
            
            Therefore, every open cover of $K$ has a finite subcover.
            Ergo, $K$ is compact.
        \end{proof}

        Next, we prove that if $K \subset (X,d)$ is compact, then every cover of $K$ consisting of open balls has a finite subcover.
        \begin{proof}
            Let $K \subset (X,d)$ be compact.
            Take any open cover of $K$ consisting of only open balls.
            Since it is an open cover of $K$, it has an open subcover of $K$.
            Therefore, every cover of $K$ consisting of open balls has a finite subcover of $K$.
        \end{proof}

    \subsection{Solution (c)}
        % \textbf{Gerry's Proposition:}
        %     Any open ball of radius $r$ can be expressed as 
        % Now we prove the initial statement.
        This uses the statement proved in part (b).
        \begin{proof}[Proof by Construction]
            Take any $\epsilon > 0$.

            Consider the collection $\mathcal{B}$ of open balls of radius $\epsilon$ defined as follows.
            For every $x \in K$, create an open ball of radius $\epsilon$ centered at $x$ ($B_\epsilon(x)$) and add that open ball to $\mathcal{B}$.
            Since every $x \in K$ is also in some open ball in $\mathcal{B}$ and every open ball is open, $\mathcal{B}$ is an open cover of $K$.

            Since $K$ is compact, $\mathcal{B}$ has a finite subcover over $K$, which we can call $\mathcal{F}$.
            Take the set of the center of every open ball $B_\epsilon(x_j)$ in $\mathcal{F}$, which will be finite because $\mathcal{F}$ is finite, and call that set $C = \{x_1, \dots, x_n\}$.
            Since the center of every open ball in $\mathcal{B}$ is in $K$ and $\mathcal{F}$ is a finite subset of $\mathcal{B}$ over $K$, the center of every open ball in $\mathcal{F}$ is in $K$.
            Therefore, $C = \{x_1, \dots, x_n\}$ is a subset of $K$.

            Since $\mathcal{F}$ is a cover of $K$, for every $y \in K$, there is an open ball of radius $\epsilon$ in $\mathcal{F}$ called $B_\epsilon(x_j)$ such that $y \in B_\epsilon(x_j)$.
            Therefore, for every $y \in K$, $d(x_j, y) < \epsilon$.
            Therefore, since the centers of all the open balls in $\mathcal{F}$ are all in $\{x_1, \dots, x_n\}$, each $y \in K$ is within distance $\epsilon$ of some $x_j$ in $\{x_1, \dots, x_n\}$.
            Ergo, for every $\epsilon > 0$, there exists a finite set $\{x_1, \dots, x_n\} \subset K$ so that each $y \in K$ is within distance $\epsilon$ of some $x_j$.
        \end{proof}
\end{document}